\chapter{Multi-Asset Options}\label{ch:dv:multi-asset-options}

A note pays a coupon at the end of the year only if all three of its shares are then at or above 70\% of their
starting levels. The client sees three familiar names and a condition that looks easy to meet. The dealer sees a
position in correlation that no listed market quotes. If the three shares were independent, with volatilities of
25\%, 30\% and 35\%, the probability of paying the coupon would be 62\%. If they moved together it would be 79\%. The
dealer who sold the coupon is short that difference. This chapter prices options on several assets, baskets,
worst-ofs and best-ofs, and shows how desks read correlation from the index and single-stock option markets. It
then covers \omterm{def:dv:multi-asset-options:dispersion}{dispersion trades} on the gap between the two, \omterm{def:dv:multi-asset-options:local}{local correlation}, which fits the index skew, and the
quanto and \omterm{def:dv:multi-asset-options:composite}{composite options} that bring a second currency into the payoff.

\section{Correlation and baskets}

Every model in this chapter starts from correlated lognormal assets:
$dS_i/S_i=\mu_i\,dt+\sigma_i\,dW^i$ with $d\langle W^i,W^j\rangle=\rho_{ij}\,dt$. Paths are simulated by the
Cholesky factor $L$ of the correlation matrix, $LL^\top=\rho$ (One Quant Book 4, chapter 25), applied to independent
normals, and prices by Monte Carlo (Book 4, chapter 26). The build takes each $\sigma_i$ as a number or as a \omterm{def:dv:local-volatility:model}{local
volatility} function $\sigma_i(t,S_i)$, so each asset's own smile is respected.

\begin{definition}[Basket option]\label{def:dv:multi-asset-options:basket}
A \emph{basket option}\index{basket option} is an option on a weighted sum of assets, typically of their
performances: $\bigl(\sum_iw_iS_{i,T}/S_{i,0}-K\bigr)^+$ for a call.
\end{definition}

A basket's variance is $\sum_{ij}w_iw_j\rho_{ij}\sigma_i\sigma_j$ in the lognormal approximation. It is below the
weighted average of the variances whenever the correlations are below one, which is the diversification the client
buys. A basket of lognormals is not lognormal, but a lognormal with the same first two moments prices it well. Match
$M_1=\sum w_iF_i$ and $M_2=\sum_{ij}w_iw_jF_iF_je^{\rho_{ij}\sigma_i\sigma_jT}$, then price with Black at the volatility
$\sqrt{\ln(M_2/M_1^2)/T}$. On an equally weighted basket of the three shares, one year, the approximation is within
0.0004 of the simulated at-the-money price for every correlation from 0 to 0.95.

\begin{omfigure}
\begin{tikzpicture}
\begin{groupplot}[group style={group size=2 by 1,horizontal sep=1.6cm},
  width=6.6cm,height=5cm,
  tick label style={font=\small},label style={font=\small},grid=major,grid style={gray!25},table/col sep=comma]
\nextgroupplot[xlabel={correlation},ylabel={price (per unit notional)},xmin=0,xmax=1,ymin=0,ymax=0.3,
  yticklabel style={/pgf/number format/fixed,/pgf/number format/precision=2},
  title={\small at-the-money options, one year},
  legend columns=1,legend cell align=left,legend style={font=\footnotesize,at={(0.5,-0.32)},anchor=north,draw=none,fill=none}]
\addplot[omThm,very thick,mark=*,mark size=1.4pt] table[x=rho,y=worst_put]{figdata/derivatives/17-multi-asset-options/three_share.csv};
\addplot[omDef,very thick,mark=square*,mark size=1.4pt] table[x=rho,y=best_call]{figdata/derivatives/17-multi-asset-options/three_share.csv};
\addplot[omProp,very thick,mark=triangle*,mark size=1.6pt] table[x=rho,y=basket_call]{figdata/derivatives/17-multi-asset-options/three_share.csv};
\legend{worst-of put,best-of call,basket call}
\nextgroupplot[xlabel={correlation},ylabel={probability},xmin=0,xmax=1,ymin=0.55,ymax=0.85,
  title={\small all three above 70\%}]
\addplot[omThm,very thick,mark=*,mark size=1.4pt] table[x=rho,y=worst_digital]{figdata/derivatives/17-multi-asset-options/three_share.csv};
\end{groupplot}
\end{tikzpicture}

\omcaption{Three shares with volatilities of 25\%, 30\% and 35\%, one year, zero rates, equal pairwise correlation.
Left: worst-of puts and best-of calls lose value as correlation rises and basket calls gain. Right: the
probability that all three end above 70\% of their start, the note's coupon condition. Data: the
tutorial.}\label{fig:dv:multi-asset-options:correlation}
\end{omfigure}

\section{Worst-of and best-of}

\begin{definition}[Worst-of option]\label{def:dv:multi-asset-options:worstof}
A \emph{worst-of option}\index{worst-of option} pays on the worst performance among several assets, for instance
$(K-\min_iS_{i,T}/S_{i,0})^+$ for a worst-of put.
\end{definition}

\begin{definition}[Best-of option]\label{def:dv:multi-asset-options:bestof}
A \emph{best-of option}\index{best-of option} pays on the best performance, for instance
$(\max_iS_{i,T}/S_{i,0}-K)^+$ for a best-of call.
\end{definition}

Correlation enters these payoffs in the opposite direction from a basket. When the assets move independently, the
worst performer is far below the average and the best far above it. When they move together, all performances are
close, and the worst-of and best-of look like options on a single asset. So a worst-of put and a best-of call lose
value as correlation rises, and a basket call gains
(\cref{fig:dv:multi-asset-options:correlation}).

\begin{example}[The note's three shares]\label{ex:dv:multi-asset-options:three}
For the three shares (volatilities 25\%, 30\% and 35\%, one year, zero rates), the at-the-money worst-of put is worth
0.250 per unit notional with independent shares and 0.152 at a correlation of 0.95. The best-of call falls from 0.284 to
0.160, and the basket call rises from 0.071 to 0.118. The probability that all three end at or above 70\% of their start
rises from 0.616 to 0.787. At a correlation of 0.4 it is 0.673, and at 0.6 it is 0.707: a five-point error in
correlation moves the coupon's value by about one point of probability.
\end{example}

A note that pays a coupon when all shares are above a level, and returns less than the capital when the worst share has
fallen through a lower level (chapter 18's autocallables), leaves the dealer long a worst-of put and short a worst-of
digital. Both lose value as correlation rises, so the dealer is short correlation, and the market for index options is
where that exposure can be partly offset. The dealer's other exposures are cross-gammas, the second derivatives in
pairs of assets (One Quant Book 6, chapter 3). They make the hedge in each share depend on the moves of the others, and
they are largest near the barrier where the worst performer changes.

\section{Implied correlation, dispersion and correlation skew}

An index is a basket of its members, so index options and single-stock options together imply a correlation.

\begin{definition}[Implied correlation]\label{def:dv:multi-asset-options:implied}
The \emph{implied correlation}\index{implied correlation} of an index with weights $w_i$ is the single correlation
$\bar\rho$ that makes the members' implied volatilities $\sigma_i$ consistent with the index's implied volatility
$\sigma_I$:
\[
\bar\rho=\frac{\sigma_I^2-\sum_iw_i^2\sigma_i^2}{\bigl(\sum_iw_i\sigma_i\bigr)^2-\sum_iw_i^2\sigma_i^2}.
\]
\end{definition}

Computed with at-the-money volatilities, it is a single number. Computed at each strike, using the index's volatility
there and the members' volatilities at the same moneyness, it varies with the strike: index skews are steeper than
single-stock skews, so the \omterm{def:dv:multi-asset-options:implied}{implied correlation} is higher at low strikes.

\begin{definition}[Correlation skew]\label{def:dv:multi-asset-options:skew}
The \emph{correlation skew}\index{correlation skew} is the variation of \omterm{def:dv:multi-asset-options:implied}{implied correlation} with the strike: higher for
index options struck below the money, where the market prices stocks falling together, than above it. It is the
equity analogue of the base-correlation skew of credit tranches (One Quant Book 2, chapter 24).
\end{definition}

\begin{example}[A synthetic 20-stock index]\label{ex:dv:multi-asset-options:index}
Take an index of twenty equally weighted stocks with flat smiles and volatilities spaced evenly from 22\% to 36\%, and let
its own one-year smile be chapter 9's: 21.9\% at 90, 19.2\% at 100, 17.0\% at 110. The \omterm{def:dv:multi-asset-options:implied}{implied correlation} is 0.549 at 90,
0.408 at 100 and 0.308 at 110 (\cref{fig:dv:multi-asset-options:skew}, left). Priced with a constant correlation of 0.408,
the members produce an index smile that is flat, 19.3\% at all three strikes: constant correlation cannot produce an index
skew out of flat member smiles. In variance-swap terms, using the index's strip volatility of 22.1\%, the \omterm{def:dv:multi-asset-options:implied}{implied correlation}
is 0.559.
\end{example}

\omterm{def:dv:multi-asset-options:implied}{Implied correlation} has been higher, on average, than the correlation stocks then realised. Driessen, Maenhout and Vilkov
report averages of 39.5\% implied against 32.5\% realised for the S\&P~500, and 46.0\% against 35.5\% for the Dow Jones
30. They read the gap as a large negative premium for correlation risk, and attribute the index \omterm{def:dv:variance-swaps-and-volatility-derivatives:vrp}{variance risk premium} of
chapter 14 to it. The trade that harvests it has a name.

\begin{definition}[Dispersion trade]\label{def:dv:multi-asset-options:dispersion}
A \emph{dispersion trade}\index{dispersion trade} sells index volatility and buys the members' volatility, in \omterm{def:dv:variance-swaps-and-volatility-derivatives:vs}{variance
swaps} or straddles, in amounts that cancel the exposure to each member's volatility. What remains is a short position in
the index's \omterm{def:dv:multi-asset-options:implied}{implied correlation}: it gains when realised correlation comes in below implied.
\end{definition}

\begin{definition}[Correlation swap]\label{def:dv:multi-asset-options:corrswap}
A \emph{correlation swap}\index{correlation swap} pays the weighted average of the pairwise correlations realised
between the members of a basket, $\sum_{i<j}w_iw_j\hat\rho_{ij}/\sum_{i<j}w_iw_j$, against a fixed strike. It isolates the
correlation that a \omterm{def:dv:multi-asset-options:dispersion}{dispersion trade} holds only approximately.
\end{definition}

The \omterm{def:dv:multi-asset-options:dispersion}{dispersion trade}'s weights follow from the index variance formula. At a fixed correlation, the derivative of the
index variance with respect to member $i$'s variance is
$w_i^2(1-\bar\rho)+\bar\rho\,w_i\sum_jw_j\sigma_j/\sigma_i$. A short index \omterm{def:dv:variance-swaps-and-volatility-derivatives:vs}{variance swap} with \omterm{def:dv:variance-swaps-and-volatility-derivatives:varnot}{variance notional} $N$,
hedged with member \omterm{def:dv:variance-swaps-and-volatility-derivatives:vs}{variance swaps} of $N$ times these derivatives, is insensitive to small changes in each member's
volatility. Its P\&L is then $N(\bar\rho-\hat\rho)C$, where $C=(\sum w_i\sigma_i)^2-\sum w_i^2\sigma_i^2$ is the
cross term of the index variance.

\begin{example}[Correlation fifteen points below implied]\label{ex:dv:multi-asset-options:dispersion}
Sell one-year index variance with a \omterm{def:dv:variance-swaps-and-volatility-derivatives:veganot}{vega notional} of 100\,000 at the strip's 22.12\%, a \omterm{def:dv:variance-swaps-and-volatility-derivatives:varnot}{variance notional} of 2\,261 per
variance point, and buy the members' variance in the hedge amounts. If each member realises its implied volatility and
correlation realises 0.409, fifteen points below the implied 0.559, the P\&L is
$2\,261\times0.15\times798=270\,600$, with $C=798$ variance points. Over 4\,000 simulated years of daily returns the
P\&L averages 271\,000, with a standard deviation of 34\,000 from the noise in realised volatilities and correlation. Nine
times in ten it lies between 215\,000 and 326\,000 (\cref{fig:dv:multi-asset-options:dispersion}).
\end{example}

\begin{omfigure}
\begin{tikzpicture}
\begin{groupplot}[group style={group size=2 by 1,horizontal sep=1.6cm},
  width=6.6cm,height=5cm,
  tick label style={font=\small},label style={font=\small},grid=major,grid style={gray!25},table/col sep=comma]
\nextgroupplot[xlabel={index strike},ylabel={implied correlation},xmin=87,xmax=113,ymin=0.25,ymax=0.6,
  xtick={90,100,110},title={\small correlation skew}]
\addplot[omThm,very thick,mark=*,mark size=2pt] table[x=k,y=implied_corr]{figdata/derivatives/17-multi-asset-options/corr_skew.csv};
\nextgroupplot[xlabel={index log-moneyness $x$},ylabel={local correlation $\rho(x)$},xmin=-0.25,xmax=0.25,ymin=0,ymax=1.05,
  title={\small calibrated local correlation}]
\addplot[omDef,very thick] table[x=x,y=rho]{figdata/derivatives/17-multi-asset-options/local_corr.csv};
\end{groupplot}
\end{tikzpicture}

\omcaption{A synthetic index of twenty stocks with flat smiles, whose own smile is chapter 9's. Left: the \omterm{def:dv:multi-asset-options:implied}{implied
correlation} at three strikes. Right: the \omterm{def:dv:multi-asset-options:local}{local correlation}, a function of the index level, that reproduces the index
smile at the three strikes. It is capped at one below the money and falls as the index rises to about 10\% above its
start; the upturn beyond is an artefact of a quadratic fitted to three strikes, which constrain nothing there. Data: the
tutorial.}\label{fig:dv:multi-asset-options:skew}
\end{omfigure}

\begin{omfigure}
\begin{tikzpicture}
\begin{axis}[width=10.5cm,height=5cm,
  xlabel={P\&L of the dispersion trade (thousands)},ylabel={share of simulations},
  tick label style={font=\small},label style={font=\small},
  xmin=100,xmax=450,ymin=0,ymax=0.15,grid=major,grid style={gray!25},
  yticklabel style={/pgf/number format/fixed,/pgf/number format/precision=2},
  table/col sep=comma]
\addplot[omThm,very thick,const plot,fill=omThm!25] table[x=pnl,y=share]{figdata/derivatives/17-multi-asset-options/dispersion.csv} \closedcycle;
\draw[black,dashed,thick] (axis cs:270.6,0) -- (axis cs:270.6,0.15);
\node[font=\footnotesize,anchor=west] at (axis cs:335,0.12) {$N(\bar\rho-\hat\rho)C=270.6$};
\end{axis}
\end{tikzpicture}

\omcaption{The P\&L of a one-year \omterm{def:dv:multi-asset-options:dispersion}{dispersion trade} (short index variance, \omterm{def:dv:variance-swaps-and-volatility-derivatives:veganot}{vega notional} 100\,000, long member variance
in the hedge amounts) over 4\,000 simulated years in which correlation realises fifteen points below implied. The dashed
line is the formula. Data: the tutorial.}\label{fig:dv:multi-asset-options:dispersion}
\end{omfigure}

Dispersion is not free money. The trade is short correlation, and the \omterm{def:dv:multi-asset-options:skew}{correlation skew} is the market's statement that
correlation rises when the index falls. In a sell-off the short index leg then loses more than the member legs gain. The limits-to-arbitrage reading of
Driessen, Maenhout and Vilkov is that realistic trading frictions eat the premium: the trade involves dozens of option
positions, rebalanced as weights and volatilities drift.

\section{Local correlation}

\cref{ex:dv:multi-asset-options:index} showed that members with flat smiles and a constant correlation cannot produce the
index's skew. The index skew is steeper than the members' skews because correlation rises when the index falls. Langnau
argued that a Gaussian copula (One Quant Book 4, chapter 15) of the single-stock distributions fails to explain the
steepness of the index skew. He extended multi-asset \omterm{def:dv:local-volatility:model}{local volatility} by making correlation a function of the market
state.

\begin{definition}[Local correlation]\label{def:dv:multi-asset-options:local}
A \emph{local correlation}\index{local correlation} model gives each asset its own \omterm{def:dv:local-volatility:model}{local volatility}, fitted to its own
options, and makes the correlation between them a function of time and of the market state, typically the index level:
$\rho_{ij}(t,I_t)$. The function is chosen so that the model's index options reprice the index smile, which it does by
construction in the same way that \omterm{def:dv:local-volatility:model}{local volatility} fits one asset's smile.
\end{definition}

Exact calibration uses the \omterm{def:dv:local-volatility:projection}{Markovian projection} of chapter 9. The index's local variance must equal the conditional
expectation of the basket's instantaneous variance given the index level, and that equation is solved for the
correlation at each index level. The build uses the simplest version that shows the mechanism. The correlation is
common to all pairs and quadratic in the index's \omterm{def:dv:implied-volatility-and-its-surface:surface}{log-moneyness} $x$, $\rho(x)=a-bx+cx^2$ clipped to $[0,1]$, and its three
coefficients are fitted to the index volatilities at the three strikes by Newton steps on a simulation with common random
numbers.

\begin{example}[Fitting the index skew]\label{ex:dv:multi-asset-options:lc}
For the synthetic index the fit gives $\rho(x)=0.384-3.62x+16.3x^2$, which reprices the index at 21.93\%, 19.19\% and
16.99\% against 21.93\%, 19.19\% and 16.98\% (\cref{fig:dv:multi-asset-options:skew}, right). Correlation is 0.91 when the
index is 10\% down in log terms, capped at one below that, and 0.19 when it is 10\% up. The members' marginal laws are
unchanged, lognormal at their own volatilities. Only their dependence moves with the market.
\end{example}

\omterm{def:dv:multi-asset-options:local}{Local correlation} then prices every product on the members consistently with both option markets. That includes worst-ofs,
baskets struck away from the money, and dispersion structures. Like \omterm{def:dv:local-volatility:model}{local volatility}, it fixes the dynamics of
correlation to a function of the state, and so it forecasts a forward \omterm{def:dv:multi-asset-options:skew}{correlation skew} that flattens, the multi-asset
counterpart of the \omterm{def:dv:local-volatility:forward}{forward smile} of chapter 16. Worst-of products whose risk sits in scenarios where the index has already
fallen inherit that assumption.

\section{Quanto and composite options}

A payoff can involve a second currency even when it references one asset. The two standard cases differ in whether the
exchange rate is fixed in advance or converted at maturity.

\begin{definition}[Quanto option]\label{def:dv:multi-asset-options:quanto}
A \emph{quanto option}\index{quanto option} pays a payoff computed on a foreign asset in the foreign currency's units but
settled in the domestic currency at a fixed exchange rate: for instance $(S_T-K)^+$ dollars per index point of a yen
index.
\end{definition}

\begin{definition}[Quanto adjustment]\label{def:dv:multi-asset-options:quantoadj}
The \emph{quanto adjustment}\index{quanto adjustment} is the change in the asset's drift when it is priced under the
domestic measure with a fixed conversion: $-\rho\sigma_S\sigma_X$, with $\sigma_X$ the volatility of the exchange rate
(domestic per foreign unit) and $\rho$ its correlation with the asset. The quanto forward is $Fe^{-\rho\sigma_S\sigma_XT}$.
\end{definition}

The adjustment follows from Girsanov's theorem (One Quant Book 4, chapter 5). Under the foreign measure the asset's drift
is the foreign carry. A domestic investor who holds the asset converted at the floating rate earns the domestic rate on
$S_tX_t$, which fixes the drift of $SX$. Dividing out the exchange rate's own drift and the covariance term
$\rho\sigma_S\sigma_X$ leaves the asset's drift under the domestic measure. The hedger of a quanto holds the foreign asset
and must keep re-hedging its foreign-currency value. When the asset and the currency are correlated, those re-hedges
systematically gain or lose, and the adjustment prices that.

\begin{example}[An index forward paid in another currency]\label{ex:dv:multi-asset-options:quanto}
A foreign index with a forward of 100, volatility 20\%, paid one year out at one domestic unit per point. With an
exchange-rate volatility of 10\% and correlations of $-0.6$, $-0.3$, $0.3$ and $0.6$, the quanto forward is 101.21, 100.60,
99.40 and 98.81. A simulation of 200\,000 paths under the domestic measure with the adjusted drift agrees to 0.005
(\cref{fig:dv:multi-asset-options:quanto}). At a correlation of $-0.3$ the at-the-money quanto call is worth 8.29 against
7.97 for the same call settled in the foreign currency.
\end{example}

\begin{definition}[Composite option]\label{def:dv:multi-asset-options:composite}
A \emph{composite option}\index{composite option} is an option on the foreign asset's value converted into the domestic
currency at the prevailing rate, $(S_TX_T-K)^+$ with a domestic strike. Its underlying $SX$ has volatility
$\sqrt{\sigma_S^2+\sigma_X^2+2\rho\sigma_S\sigma_X}$.
\end{definition}

A quanto removes the currency risk from the payoff and leaves the correlation in the drift. A composite keeps the currency
risk and puts the correlation in the volatility: at $\rho=-0.3$, the composite volatility is $\sqrt{0.04+0.01-0.012}=19.5\%$.
Both are two-asset options, and both carry the same correlation position, which the desk must mark and hedge. Spread options
between two assets, the other common two-asset product, appear in One Quant Book 6, chapter 16.

\begin{omfigure}
\begin{tikzpicture}
\begin{axis}[width=10.5cm,height=5cm,
  xlabel={correlation between the index and the exchange rate},ylabel={quanto forward},
  tick label style={font=\small},label style={font=\small},
  xmin=-0.7,xmax=0.7,ymin=98.5,ymax=101.5,grid=major,grid style={gray!25},
  legend columns=2,legend style={font=\footnotesize,at={(0.5,-0.3)},anchor=north,draw=none,fill=none,/tikz/every even column/.append style={column sep=8pt}},
  table/col sep=comma]
\addplot[omThm,very thick] table[x=rho,y=analytic]{figdata/derivatives/17-multi-asset-options/quanto.csv};
\addplot[black,only marks,mark=*,mark size=1.8pt] table[x=rho,y=mc]{figdata/derivatives/17-multi-asset-options/quanto.csv};
\legend{$Fe^{-\rho\sigma_S\sigma_XT}$,simulation}
\end{axis}
\end{tikzpicture}

\omcaption{The one-year forward of a foreign index (forward 100, 20\% volatility) paid at a fixed exchange rate whose
volatility is 10\%, by the correlation between the index and the rate. Data: the tutorial.}\label{fig:dv:multi-asset-options:quanto}
\end{omfigure}

\section{Tutorial: correlation in practice}\label{tut:dv:multi-asset-options:tut}

\begin{tutorial}
\textbf{Goal.} Price baskets, worst-ofs and best-ofs against correlation; read the \omterm{def:dv:multi-asset-options:implied}{implied correlation} of an index at three
strikes; fit a \omterm{def:dv:multi-asset-options:local}{local correlation} to the index skew; simulate a \omterm{def:dv:multi-asset-options:dispersion}{dispersion trade}; measure the \omterm{def:dv:multi-asset-options:quantoadj}{quanto adjustment}.
\textbf{End state:} the four figures and the numbers of the weekend problem.
\begin{tutsteps}
\item \textbf{\omterm{def:dv:multi-asset-options:local}{Local correlation} paths.} Equicorrelated shocks whose common correlation is read from the index level at every
step:
\omcode{code/firm/multiasset/firm_multiasset.py}{51}{75}{The local-correlation path generator.}\label{lst:dv:multi-asset-options:lc}
\item \textbf{\omterm{def:dv:multi-asset-options:implied}{Implied correlation} and the index variance}, the two directions of the same formula:
\omcode{code/firm/multiasset/firm_multiasset.py}{119}{130}{Implied correlation from index and member volatilities.}\label{lst:dv:multi-asset-options:implied}
\item \textbf{The quanto forward and the composite volatility}:
\omcode{code/firm/multiasset/firm_multiasset.py}{143}{151}{Quanto and composite adjustments.}\label{lst:dv:multi-asset-options:quanto}
\item \textbf{Run} \texttt{dv\_multiasset.three\_share\_prices()}, \texttt{correlation\_by\_strike()},
\texttt{calibrate\_local\_correlation()}, \texttt{dispersion()}, \texttt{quanto\_example()} and \texttt{fig\_multiasset.py}.
\end{tutsteps}
\textbf{What to change next.} Give the members their own skews (a \omterm{def:dv:local-volatility:model}{local volatility} each) and refit the \omterm{def:dv:multi-asset-options:local}{local correlation};
price the three-share worst-of digital under the fitted \omterm{def:dv:multi-asset-options:local}{local correlation}; run the \omterm{def:dv:multi-asset-options:dispersion}{dispersion trade} with correlation realising
at implied but member volatilities five points below.
\end{tutorial}

\section{Build: multi-asset paths and payoffs}\label{bld:dv:multi-asset-options:multiasset}

\begin{build}
\bfield{Purpose} The miniature firm's multi-asset engine: correlated paths for baskets, worst-ofs and quantos, the implied-
and realised-correlation analytics of its dispersion book, and the generator that chapter 18's autocallables run on.
\bfield{Interface} \texttt{correlated\_paths(s0, vols, corr, times, r, q, n\_paths, seed, quanto, steps\_per\_date)} with
\texttt{vols[i]} a number or a function $\sigma_i(t,S)$; \texttt{local\_correlation\_paths(s0, vols, weights, rho\_fn, times,
\ldots)}; \texttt{equicorrelation}; \texttt{basket\_payoff}, \texttt{worst\_of\_payoff}, \texttt{best\_of\_payoff},
\texttt{worst\_of\_digital}; \texttt{basket\_moment\_match}; \texttt{implied\_correlation}, \texttt{index\_variance},
\texttt{realised\_correlation}; \texttt{quanto\_forward}, \texttt{composite\_vol}.
\bfield{Rules} Correlation matrices are checked positive definite before use (the Cholesky factor fails otherwise); the quanto
correlation is quoted against the rate in domestic per foreign units; every Monte Carlo number carries a standard error.
\bfield{Acceptance tests} \texttt{code/firm/multiasset/tests/}: simulated assets are martingales with the specified correlation;
worst-of and best-of reduce to the vanilla at correlation one; moment matching against simulation; \omterm{def:dv:multi-asset-options:implied}{implied correlation} inverts
the index variance; \omterm{def:dv:multi-asset-options:local}{local correlation} reduces to the constant case; the quanto forward matches a simulation.
\bfield{Stretch} Per-member local volatilities and a pairwise \omterm{def:dv:multi-asset-options:local}{local correlation} calibrated by \omterm{def:dv:local-volatility:projection}{Markovian projection}; the
cross-gamma report of One Quant Book 6.
\end{build}

\begin{omsources}
\item J.~Driessen, P.~Maenhout and G.~Vilkov, ``The price of correlation risk: evidence from equity options'', SSRN 673425
(2005); ``Option-implied correlations and the price of correlation risk'', SSRN 2359380 (2013).
\item A.~Langnau, ``Introduction into `local correlation modelling''', arXiv 0909.3441 (2009).
\end{omsources}

\section{Exercises}

\begin{exercise}[$\star$]\label{exo:dv:multi-asset-options:1}
Two assets with volatilities 20\% and 30\%, equally weighted, correlation 0.5: what is the basket's volatility in the lognormal
approximation?
\end{exercise}

\begin{exercise}[$\star$]\label{exo:dv:multi-asset-options:2}
Why does a best-of call lose value as correlation rises while a basket call gains?
\end{exercise}

\begin{exercise}[$\star$]\label{exo:dv:multi-asset-options:3}
An index of two equally weighted members with volatilities 25\% and 35\% trades at 24\%. Compute its \omterm{def:dv:multi-asset-options:implied}{implied correlation}.
\end{exercise}

\begin{exercise}[$\star\star$]\label{exo:dv:multi-asset-options:4}
Derive the quanto forward $Fe^{-\rho\sigma_S\sigma_XT}$ from the requirement that $S_TX_T$, converted at the floating rate,
earns the domestic rate.
\end{exercise}

\begin{exercise}[$\star\star$]\label{exo:dv:multi-asset-options:5}
Explain why a dealer who sells worst-of notes is short correlation, and name one listed instrument that offsets part of the
position.
\end{exercise}

\begin{exercise}[$\star\star$]\label{exo:dv:multi-asset-options:6}
Using the synthetic index, compute the dispersion P\&L if realised correlation comes in at implied but every member's volatility
realises one point above its implied level (hedge amounts unchanged).
\end{exercise}

\begin{exercise}[$\star\star\star$]\label{exo:dv:multi-asset-options:7}
\emph{Coding.} Price the three-share worst-of digital (all above 70\%) under the synthetic index's fitted \omterm{def:dv:multi-asset-options:local}{local correlation}
function, applied to the three shares' own basket. Compare with the constant-correlation prices of the chapter.
\end{exercise}

\begin{exercise}[$\star\star\star$]\label{exo:dv:multi-asset-options:8}
\emph{Find the flaw.} ``\omterm{def:dv:multi-asset-options:implied}{Implied correlation} has exceeded realised correlation on average, so a dispersion book that sells index
variance and buys member variance is a riskless carry trade.''
\end{exercise}

\section{Problem: The Index and Its Members}

\begin{problem}[{Weekend problem --- implied correlation and a dispersion trade}]\label{pb:dv:multi-asset-options:1}
A synthetic index holds twenty stocks in equal weights, with flat smiles and volatilities evenly spaced from 22\% to 36\%. The
index's own one-year smile is chapter 9's. A dispersion desk considers selling index variance against member variance.

\medskip\noindent\textbf{Part I --- The members.}
\begin{enumerate}
\item Compute $\sum w_i^2\sigma_i^2$ and the cross term $C$.
\item What index volatility would a correlation of zero imply? Of one?
\item Why is the index volatility below the average member volatility?
\item Give the index's volatility at 90, 100 and 110.
\item Give the index volatility that a constant correlation of 0.408 produces at 90 and 110.
\end{enumerate}

\medskip\noindent\textbf{Part II --- \omterm{def:dv:multi-asset-options:implied}{Implied correlation}.}
\begin{enumerate}[resume]
\item Give the \omterm{def:dv:multi-asset-options:implied}{implied correlation} at 90, 100 and 110.
\item Give the variance-swap \omterm{def:dv:multi-asset-options:implied}{implied correlation}.
\item Explain the \omterm{def:dv:multi-asset-options:skew}{correlation skew} in one sentence.
\item Give the \omterm{def:dv:multi-asset-options:local}{local correlation} function that reprices the three strikes.
\item What \omterm{def:dv:multi-asset-options:local}{local correlation} does it assign 10\% down and 10\% up?
\end{enumerate}

\medskip\noindent\textbf{Part III --- The \omterm{def:dv:multi-asset-options:dispersion}{dispersion trade}.}
\begin{enumerate}[resume]
\item For a \omterm{def:dv:variance-swaps-and-volatility-derivatives:veganot}{vega notional} of 100\,000 on the index, give the \omterm{def:dv:variance-swaps-and-volatility-derivatives:varnot}{variance notional}.
\item Give the member hedge amounts' formula and explain what they cancel.
\item Give the P\&L if correlation realises fifteen points below implied.
\item Give the simulated mean, standard deviation and 5--95\% range.
\item What would a \omterm{def:dv:multi-asset-options:corrswap}{correlation swap} struck at the implied 0.559 pay per unit notional?
\end{enumerate}

\medskip\noindent\textbf{Part IV --- Judgement.}
\begin{enumerate}[resume]
\item When does the \omterm{def:dv:multi-asset-options:dispersion}{dispersion trade} lose, and how much could it lose?
\item What does the published evidence say about implied against realised correlation?
\item Why might the premium persist?
\item State the \emph{named result}: the \omterm{def:dv:multi-asset-options:implied}{implied correlation} of the synthetic index at three strikes, and the P\&L of the
\omterm{def:dv:multi-asset-options:dispersion}{dispersion trade} when realised correlation comes in fifteen points below implied.
\item In one sentence: what does an index option price that its members' options do not?
\end{enumerate}
\end{problem}

\section{Interview questions}

\begin{interviewq}[$\star$ \iqroles{trader, researcher}]\label{iq:dv:multi-asset-options:1}
How does correlation affect the prices of a basket call, a worst-of put and a best-of call?
\end{interviewq}

\begin{interviewq}[$\star$ \iqroles{researcher}]\label{iq:dv:multi-asset-options:2}
Define \omterm{def:dv:multi-asset-options:implied}{implied correlation}. Why is it higher for low strikes?
\end{interviewq}

\begin{interviewq}[$\star\star$ \iqroles{trader}]\label{iq:dv:multi-asset-options:3}
Explain a \omterm{def:dv:multi-asset-options:dispersion}{dispersion trade}: the legs, the weights and what you are exposed to.
\end{interviewq}

\begin{interviewq}[$\star\star$ \iqroles{researcher, developer}]\label{iq:dv:multi-asset-options:4}
Derive the \omterm{def:dv:multi-asset-options:quantoadj}{quanto adjustment}. What is the sign for a Japanese index paid in dollars if the index tends to fall when the yen
strengthens?
\end{interviewq}

\begin{interviewq}[$\star\star$ \iqroles{developer}]\label{iq:dv:multi-asset-options:5}
Your correlation matrix for 30 assets fails the Cholesky factorisation. Why, and what do you do?
\end{interviewq}

\begin{interviewq}[$\star\star\star$ \iqroles{risk, trader}]\label{iq:dv:multi-asset-options:6}
A book of worst-of autocallables is short correlation. How would you measure, hedge and reserve for that exposure?
\end{interviewq}
